% Auto-generated — do not edit by hand.
\begin{table}[t]
\centering
\caption{Extended baseline comparison on the IIWA14 combined scene.}
\label{tab:baselines_v2}
\small
\begin{tabular}{@{}lrrrr@{}}
\toprule
\textbf{Method} & \textbf{Precomp.\,(s)} & \textbf{Query\,(s)}
                & \textbf{Path\,(rad)} & \textbf{SR\,(\%)} \\
\midrule
SBF (ours)              & \textbf{2.23} & 0.279 & 2.865 & \textbf{100} \\
\quad IQR               & $[1.61,2.49]$ & --    & --    & --  \\
SBF~+~GCS post (ours)   & 2.23          & 0.064 & \textbf{2.063} & \textbf{100} \\
\midrule
IRIS-NP~+~GCS           & 602           & 0.351 & 2.631  & 100 \\
IRIS-ZO~+~GCS           & 16.0          & 0.231 & --     & 40  \\
RRT-Connect             & 0             & 0.157 & 5.234  & 100 \\
PRM ($N{=}10\mathrm{k}$)& 3.23         & \textbf{0.042} & 3.703 & 100$^{\dagger}$ \\
\bottomrule
\end{tabular}
\\[2pt]
{\footnotesize $^{\dagger}$The PRM SR=$100\%$ entry comes from a
$5$-seed $\times$ $5$-query draw at $N{=}10{,}000$ on the canonical
benchmark. The 3-seed sweep of \Cref{tab:prm_sweep} shows
$80\%$ at $N{=}10{,}000$ and $93\%$ at $N{=}30{,}000$ on a
different (and broader) seed/query distribution; the two protocols
are not directly comparable and should not be blended into a single
asymptote. The IRIS-ZO row uses Drake-default probabilistic
parameters and is not the result of a parameter sweep, so its
SR=$40\%$ should be read as a single-point reference rather than as
the asymptotic IRIS-ZO behaviour.}
\end{table}
